\subsection{The characteristic-prime fibre and local propagation}
\label{sec:selmer-integral-frontier}

The integral argument requires the weight-zero subring of the full
normalized global deformation ring to be finite over the coefficient
discrete valuation ring. Rigidity after inverting its uniformizer does
not imply this finiteness, even for a reduced ring.

Indeed, let $\mathcal O$ have uniformizer $\pi$ and residue field $k$.
For $M\geq1$, give $x,y$ weights $1,-1$ and put
\[
 S_M=\mathcal O[[x,y]]/(\pi^M x,\pi^M y).
\]
Then $S_M[1/\pi]=\operatorname{Frac}(\mathcal O)$, whereas its completed
weight-zero subring is $\mathcal O[[z]]/(\pi^M z)$, with $z=xy$.
Its reduction is $k[[z]]$, so it is not finite over $\mathcal O$.
For $M=1$, the identity
$(\pi x,\pi y)=(\pi)\cap(x,y)$ shows that $S_1$ is reduced.
Passing to the $\pi$-torsion-free quotient gives $\mathcal O$ but
discards the negative augmentation conormal module
$(\mathcal O/\pi^M)y$. The curve $x=y=t$, $\pi=0$ into $k[[t]]$
also shows directly that $xy\notin\sqrt{\pi S_M}$.

These examples are not asserted to arise from Galois representations.
They show why an arithmetic assertion about the characteristic-prime
fibre is necessary. For a weight-zero subring with finitely many
topological generators, the relevant criterion places every nonconstant
generator in $\sqrt{\pi S}$. Contraction in the original grading and
topological Nakayama then give the desired finite module. In the deformation
problem this criterion is supplied by propagation through chains of
reducible, nonsemisimple representations. The following local lemma
proves the prime inclusion used in that propagation. Its application
retains the prescribed deformation relations and the dimension of the
patched automorphic module.
